% ============================ 六、开放贡献与长期研究潜力 ============================
\section{开放贡献与长期研究潜力}
\label{sec:oss}

修订说明第四条新设「开源贡献」维度（占 5\,\%），并明确要求在提交物中说明六项内容：\textbf{开源或开放的范围、开源协议或授权方式、第三方依赖、商业 API 调用情况、闭源模型使用情况、数据来源与授权边界}。本节逐项作答，然后给出两个可直接复用组件、一个可审计构建流程与长期研究计划。

\subsection{开源范围与协议}

项目仓库以 \textbf{Apache License 2.0} 发布：\\
{\small\url{https://github.com/methylation5mc2026-jpg/GOAI-proteomics-perturbtion-response-prediction}}

\begin{table}[H]
\centering\footnotesize
\caption{\textbf{开源清单与边界。} 「开放」表示随代码库发布；「受限」表示因数据使用协议或体量原因不随库发布，但提供可复现的生成脚本。}
\label{tab:oss}
\begin{tabularx}{\textwidth}{@{}R{5.5cm}cX@{}}
\toprule
内容 & 状态 & 说明与许可 \\
\midrule
全部 73 个科学工作流脚本（25\,591 行）与共享路径模块 & 开放 & 按阶段编号；主要随机训练阶段固定种子 42；Apache-2.0 \\
评分 harness 与 57 项自检 & 开放 & 纯函数模块，\textbf{可独立用于校验任何方案的自报分数}；Apache-2.0 \\
快速代理评分器与等价性门 & 开放 & 含五个对抗性校验用例；Apache-2.0 \\
测试队列自评脚本（本次新增） & 开放 & 含泄漏保护与「仅评分不选择」协议实现；Apache-2.0 \\
自演化环路运行器与日志模式 & 开放 & 含变异算子定义、Pareto 保留规则与诊断映射；Apache-2.0 \\
生物先验构建脚本（ITPV / STRING / iMM904 / 复合物） & 开放 & 含药理学正负对照清单与验证脚本；Apache-2.0 \\
外部资源披露清单生成器 & 开放 & 从产物生成机器可读披露 JSON/CSV；Apache-2.0 \\
公开结构化证据与完整性索引 & 开放 & 70 份 JSON/CSV 快照加 1 份 SHA-256 索引；不含大型或受限产物 \\
\midrule
官方竞赛数据集 & \textbf{受限} & 受赛事数据使用协议约束，\textbf{不再分发}；仓库仅含读取接口与结构契约检查 \\
完整 ITPV、ChEMBL 活性行与靶点映射表 & \textbf{未公开} & 仓库仅含构建/验证代码与统计摘要；再生成和再分发须遵守 ChEMBL CC BY-SA 3.0 \\
完整运行日志 & \textbf{未公开} & 报告只保留可由公开代码与快照交叉核对的 9 项代表性缺陷 \\
中间缓存（折外预测矩阵等，约 2.5\,GB） & \textbf{受限} & 体量原因；提供一键重建脚本 \\
测试集预测 CSV（约 199\,MB / 阶段） & \textbf{受限} & 按赛事要求提交，不公开分发 \\
模型权重（GBDT/CatBoost/PyTorch 检查点） & \textbf{受限} & 体量原因（约 1.6\,GB）；训练脚本与种子已开放，可重建 \\
\bottomrule
\end{tabularx}
\end{table}

选择 Apache-2.0 而非 MIT 的理由是它含明确的专利授权与商标条款。该许可只覆盖本仓库自有代码与文本，\textbf{不会覆盖或改写第三方数据条款}。公开仓库未发布完整 ChEMBL 衍生矩阵；若下游重新生成并再分发，需按 ChEMBL 的 \textbf{CC BY-SA 3.0} 处理相应衍生产物。

\subsection{第三方依赖、商业 API 与闭源模型}

表 \ref{tab:deps} 只列主要直接依赖，完整锁定图以 \file{uv.lock} 为准；本仓库不再作「全部传递依赖均为宽松许可」这一未经逐包归档的概括。发布者在再分发环境时仍应做独立许可证审计。外部数据侧，ChEMBL 的 CC BY-SA 3.0 边界已单列；复合物先验是脚本中的手工名册，并非随库发布的 CORUM / Complex Portal 数据导出，因此不把上游数据库许可误写成本仓库文件的既成许可。

\begin{table}[H]
\centering\small
\caption{\textbf{商业服务与闭源模型的使用披露。} 判定标准很简单：\textbf{若移除该服务，本方案报告的分数是否改变}。}
\label{tab:disclosure}
\begin{tabularx}{\textwidth}{@{}R{3.3cm}R{4.0cm}R{2.1cm}X@{}}
\toprule
资源 & 使用环节 & 是否影响分数 & 替代方案与迁移成本 \\
\midrule
大语言模型辅助的编码代理与文献核验代理 & 脚本起草与重构、文献检索与 DOI 核验、图件生成、排版 & \textbf{否}（不参与训练、推理或评分） & 全部产物为确定性脚本，可由人工重写；迁移成本为工时，不影响数值结果 \\
AI 图件生成服务 & 若干\emph{概念示意}图（含图 \ref{fig:ga}、\ref{fig:loop}） & \textbf{否} & 承载数值结论的图件全部由 Matplotlib 从产物直接绘制，可完全复现 \\
ChEMBL / STRING / UniProt / PubChem Web 服务 & 实体表征构建（一次性抽取） & \textbf{是}（构成 ITPV、图结构与分子特征） & 公开仓库仅含统计摘要，未含完整抽取表与缓存；重新建模需要联网或另行提供已冻结缓存 \\
在线推理服务（用于模型本身） & \textbf{无} & 否 & 全部训练、推理与评分均在本地完成 \\
闭源模型（作为方案组件） & \textbf{无} & 否 & 本方案不含任何闭源模型组件 \\
\bottomrule
\end{tabularx}
\end{table}

\textbf{辅助工具使用边界（按公开产物可核查）。} 编码、文献检索、概念图与排版过程使用过自动化辅助工具，但没有在线推理服务或闭源模型作为训练、推理或评分组件。当前公开仓库不含完整过程日志，因此不对每一次撰写或决策动作作无法外部复核的强断言。可核查的边界是：报告分数来自结构化快照与本地评分代码；所有承载数值结论的图件由脚本从产物绘制；概念示意不作为实验结果证据。

\subsection{两个可直接复用组件与一个可审计构建流程}

我们不认为「把代码扔上 GitHub」等于开源贡献。真正的判据是：\emph{别人能否在明确输入与许可边界下独立核查或复用}。评分 harness 与快速代理评分器可直接复用；ITPV 部分当前开放的是构建/验证流程与统计摘要，不把未发布的完整矩阵说成现成开放数据。

\begin{enumerate}
\item \textbf{评分 harness（\file{harness.py} $+$ 57 项自检）}。纯函数、无副作用，输入 $(\bm{y},\hat{\bm{y}},\bm{\Delta},\hat{\bm{\Delta}},\text{meta},\bm{\mu}_{\text{ctx}},\bm{\mu}_{\text{drug}})$，输出七模块分与四种聚合口径。\textbf{任何参赛队都可以用它核对自己的自报分数}，也可以用它检验我们的。它还内置了对平凡解的诊断（空值下限、批次感知残差敏感性），这部分对任何采用「减去群体均值再算相关」口径的扰动预测基准都直接可用。
\item \textbf{快速代理评分器与等价性门}。在权重是仿射参数的情形下把评分从 $\sim$20\,s 降到 $\sim$4.4\,ms（约 6\,184 倍），并附一道\emph{不可绕过}的五用例等价性校验。这套「加速 $+$ 强制等价性门」的模式可以迁移到任何以掩码相关系数为目标的堆叠搜索。
\item \textbf{ITPV 构建与验证流程}。代码把 ChEMBL 的 13\,963 条可用效力记录，经 UniProt 返回的 OrthoDB 交叉引用投射到 5\,243 个被测酵母蛋白坐标轴，并附事先声明的药理学正负对照。公开快照可核对覆盖率与对照结果；完整活性表、映射表和矩阵未随库发布，重新构建需要外部服务并受其许可约束。它在本任务上未带来增益，这一负结果仍保留。
\end{enumerate}

\subsection{长期研究潜力：预注册的下一阶段工作}

\begin{table}[H]
\centering\footnotesize
\renewcommand{\arraystretch}{1.05}
\caption{\textbf{下一阶段五项预注册工作及其成功判据。} 判据在实施\emph{之前}写下，以避免报告数字被回溯拟合到目标上；未达成即按 ITPV 的先例记为「可验证但无增益」。工作项 E 为本次修订新增，且被列为第一优先。}
\label{tab:plan}
\begin{tabularx}{\textwidth}{@{}p{0.45cm}R{3.2cm}XR{4.3cm}@{}}
\toprule
& \textbf{工作项} & \textbf{内容} & \textbf{预注册成功判据} \\
\midrule
\textbf{E} & \textbf{菌株基因组表征}\newline\NEW{}\textbf{（第一优先）} & 用 1011 酵母基因组项目\citep{peter2018_1011genomes}的变异集为 6 个菌株构造表征（基因存在/缺失、拷贝数、非同义变异负荷、通路层聚合），DHY210 以 SGD S288C 参考基因组\citep{engel2022sgd}作代理；作为实体特征进入两轨 & S2 模块提升 $\ge0.010$，\textbf{且 embedding 打乱对照下增益必须消失}；若打乱后分数不掉，则判定为「涨的不是基因组信息」并如实报告 \\
A & 删失似然预测头 & 实现式 \eqref{eq:tobit} 的 Tobit 头；$\tau_p$ 以逐蛋白观测最小值低分位数初始化并允许按仪器偏移，$\sigma_p$ 与预测头联合学习 & 折外队列上 Module 1 的逐蛋白 $R^2$ 提升 $\ge0.01$，\emph{或} Module 4 的 AUPRC 提升 $\ge0.01$；消融须在其他一切不变下进行 \\
B & 转录调控网络先验 & 引入 TF-靶基因边\citep{Harbison2004RegulatoryCode,Teixeira2018YEASTRACT,Kemmeren2014RegulatoryNetworks}构造调控子层低秩响应基，以与复合物一致性同构的形式加入式 \eqref{eq:mechloss} & S1/S2 残差模块提升 $\ge0.005$ \\
C & 不确定度感知的 M1 收复 & 以 Tobit 的删失概率做逐单元收缩，针对「零假设锚点在 M1 上胜过我们」这一已知缺口；把蛋白簇索引改由 GNN 中间表征给出（仅用折外拟合行生成，保留静态簇为对照） & M1 模块 $\ge$ 对照锚点水平（测试队列 0.8484）\textbf{且} M2/S1/S2 不下降 \\
D & 逐先验消融与失败模式报告 & 去 ITPV / 去 STRING 图 / 去复合物罚 / 去通量罚 / 去交叉注意力 / 去分层收缩，各自的折外与两个队列分数 & 作为提交\textbf{必需项}而非可选项 \\
\bottomrule
\end{tabularx}
\end{table}

\textbf{被明确排除的三条路线及理由。}（一）不对缺失做任何填补以「提高数据量」——这会把左删失偏倚引入训练目标；（二）不在验证集上选择任何超参数或权重，即使这会使报告分数低于可达上界（阶段 4 的验证调优上界 0.5054、阶段 5 的 0.5192 均已作为\textbf{乐观上界}标注而非成绩）；（三）不追求单纯提高 S1 的批次拟合技巧——我们已在两个队列上证明该模块的基线优势来自批次复现，把它作为优化目标等于\emph{优化孔板偏移并称之为生物学}。

\textbf{算力安排。} 当前神经成员的收敛在时间预算内被截断（阶段 6 最细簇层坐标上升带 \code{truncated} 标志）。后续计划以云端多卡实例扩大轮数与贝叶斯超参搜索；但这是\textbf{扩展项而非救命项}——负结果三已实测「150 $\to$ 300 轮」的收益为 $+0.0108$ dev-PCC / 0.547 GPU-小时，成本线性而收益递减。相较之下，工作项 E（菌株表征）不需要更多算力，只需要正确的外部数据——这也是我们把它排在第一优先的原因。

\begin{keybox}[title=结语]
\footnotesize
本方案不只报告验证分数从 0.445 到 0.509、测试队列自评 0.464，更把题目组织成\textbf{可被质询的环境}：零假设在两个队列上核算，评分漏洞被独立复现，知识注入受事先声明的对照检验，未达成目标、失效架构与\textbf{尚未补齐的菌株表征}均明确披露。统一开放榜要求也使未见实体表征成为下一阶段第一优先；公开 harness 与检查门用于后续独立核查。
\end{keybox}
